%%%%%%%%%%%%%%%%%%%%%%%%%%%%%%%%%%%%%%%%%%%%%%%%%%%%%%%%%%%%%%%%%%%%%%%%%
%% Linear Algebra -- course project report template, Autumn 2026
%% Ukrainian Catholic University, Faculty of Applied Sciences,
%% Data Science Master Programme
%%
%% The title page is adapted from a template by WikiBooks, distributed
%% by LaTeXTemplates.com under CC BY-NC-SA 3.0.
%%
%% HOW TO USE THIS TEMPLATE
%%   1. Fill in the block marked "FILL IN" below: the stage of the report,
%%      the title, the authors, the project card and the links.
%%   2. One file serves all three milestones. Keep working in the same
%%      file: interim report 1 grows into interim report 2, and that grows
%%      into the final report.
%%   3. Replace the grey guidance boxes with your own text. Before you
%%      submit, set \showguidefalse so that any box you forgot disappears.
%%   4. Do NOT change the font size, the paper size or the margins.
%%      Page limits are counted in this format.
%%   5. Compile with pdfLaTeX (the Overleaf default). The folder must
%%      contain UCU-Apps.png (the logo); the example figure uses an image
%%      shipped with every TeX distribution.
%%
%% PAGE LIMITS (counted automatically; see the end of the log file)
%%   Interim reports 1 and 2: 2 pages, project card included.
%%   Final report: at most 10 pages from the abstract to the end of the
%%   Conclusion, figures and tables included. The title page, the
%%   references and the appendices are not counted.
%%%%%%%%%%%%%%%%%%%%%%%%%%%%%%%%%%%%%%%%%%%%%%%%%%%%%%%%%%%%%%%%%%%%%%%%%

\documentclass[11pt,a4paper]{article}

%----------------------------------------------------------------------
%  Page layout: do not change
%----------------------------------------------------------------------
\usepackage[a4paper,left=20mm,right=20mm,top=20mm,bottom=25mm,
            headheight=14pt,marginparwidth=16mm]{geometry}

%----------------------------------------------------------------------
%  Packages
%----------------------------------------------------------------------
\usepackage[english]{babel}
% T1 encoding only together with the scalable Latin Modern fonts
% (standard on Overleaf); otherwise fall back to the default encoding.
\IfFileExists{lmodern.sty}{\usepackage[T1]{fontenc}\usepackage{lmodern}}{}
\usepackage{microtype}
\usepackage{amsmath,amssymb,amsthm}
\usepackage{graphicx}
\usepackage{booktabs,tabularx}
\usepackage{caption,subcaption}
\usepackage[dvipsnames]{xcolor}
\usepackage{fancyhdr}
\usepackage{refcount}
\usepackage{listings}
\usepackage[ruled,vlined,linesnumbered]{algorithm2e}
% Margin notes are useful while drafting. Add the option "disable"
% before you submit, or simply delete your \todo{} commands.
\usepackage[colorinlistoftodos,textsize=scriptsize]{todonotes}
\usepackage[colorlinks=true,linkcolor=MidnightBlue,citecolor=MidnightBlue,
            urlcolor=MidnightBlue]{hyperref}   % load hyperref last

%%%%%%%%%%%%%%%%%%%%%%%%%%%%%%%%%%%%%%%%%%%%%%%%%%%%%%%%%%%%%%%%%%%%%%%%%
%%  FILL IN
%%%%%%%%%%%%%%%%%%%%%%%%%%%%%%%%%%%%%%%%%%%%%%%%%%%%%%%%%%%%%%%%%%%%%%%%%

% Stage of the report: 1 = interim report 1 (scoping memo),
%                      2 = interim report 2 (working baseline),
%                      3 = final report
\newcommand{\stage}{1}

\newcommand{\reporttitle}{Low-Rank Spectral Fingerprints for AI-Generated Music Detection}
\newcommand{\shorttitle}{Low-rank spectrograms for AI-music detection}
\newcommand{\reportauthors}{Iryna Matvieieva \quad Petro Mozil}

% Project card: one line each. Keep it consistent with your registration;
% if something has changed, update it here and say so in the report.
\newcommand{\carddomain}{Audio forensics / music provenance}
\newcommand{\cardmethod}{STFT magnitude spectrograms; exact and randomized SVD}
\newcommand{\cardobject}{For each 10-s mono clip, $M\in\mathbb{R}_{\geq0}^{1025\times T}$, with $T\approx431$; analyse the leading $k=20$ singular values}
\newcommand{\cardalgorithms}{Own randomized SVD ($k=20$, oversampling $p=10$, one power step); exact NumPy/LAPACK SVD as reference}
\newcommand{\carddata}{FakeMusicCaps (subset), paired with retrievable original MusicCaps clips~\cite{FakeMusicCaps}}
\newcommand{\cardcode}{\url{https://github.com/pmozil/music-ai-spectra}}
\newcommand{\cardvideo}{\url{https://link-to-your-recording}} % final only

% Guidance boxes: set \showguidefalse before you submit.
\newif\ifshowguide
\showguidetrue

%%%%%%%%%%%%%%%%%%%%%%%%%%%%%%%%%%%%%%%%%%%%%%%%%%%%%%%%%%%%%%%%%%%%%%%%%
%%  Template machinery: no need to read further than the next block
%%%%%%%%%%%%%%%%%%%%%%%%%%%%%%%%%%%%%%%%%%%%%%%%%%%%%%%%%%%%%%%%%%%%%%%%%

% Notation used in the course. Add your own macros here.
\newcommand{\bR}{\mathbb{R}}
\newcommand{\bC}{\mathbb{C}}
\newcommand{\ba}{\mathbf{a}}
\newcommand{\bb}{\mathbf{b}}
\newcommand{\bu}{\mathbf{u}}
\newcommand{\bv}{\mathbf{v}}
\newcommand{\bx}{\mathbf{x}}
\newcommand{\by}{\mathbf{y}}
\newcommand{\bz}{\mathbf{z}}
\DeclareMathOperator{\rank}{rank}
\DeclareMathOperator{\cond}{cond}
\DeclareMathOperator{\diag}{diag}
\DeclareMathOperator{\trace}{trace}
\DeclareMathOperator{\nnz}{nnz}
\DeclareMathOperator*{\argmin}{argmin}

% Theorem-like environments, numbered within sections
\theoremstyle{plain}
\newtheorem{theorem}{Theorem}[section]
\newtheorem{proposition}[theorem]{Proposition}
\newtheorem{lemma}[theorem]{Lemma}
\theoremstyle{definition}
\newtheorem{definition}[theorem]{Definition}
\newtheorem{example}[theorem]{Example}
\theoremstyle{remark}
\newtheorem{remark}[theorem]{Remark}

% Code listings
\lstset{language=Python, basicstyle=\ttfamily\small, columns=fullflexible,
  keepspaces=true, frame=lines, numbers=left, numberstyle=\tiny\color{gray},
  commentstyle=\color{gray}, keywordstyle=\color{MidnightBlue}\bfseries,
  showstringspaces=false, breaklines=true}

% Guidance boxes
\newcommand{\guide}[1]{\ifshowguide\par\smallskip\noindent
  \colorbox{gray!12}{\parbox{\dimexpr\linewidth-2\fboxsep}{\small\itshape
  \textbf{Guidance.} #1}}\par\smallskip\fi}

% Stage names
\newcommand{\stagename}{\ifcase\stage\or Interim report 1: scoping memo\or
  Interim report 2: working baseline\or Final report\fi}

% Running heads
\pagestyle{fancy}
\fancyhf{}
\fancyhead[L]{\small\shorttitle}
\fancyhead[R]{\small\stagename}
\fancyfoot[C]{\small\thepage}
\renewcommand{\headrulewidth}{0.4pt}

% Project card
\newcommand{\projectcard}{%
  \noindent
  \begin{tabularx}{\linewidth}{@{}>{\raggedright\bfseries}p{3.0cm}
                                  >{\raggedright\arraybackslash}X@{}}
    \toprule
    Domain        & \carddomain \\
    Method family & \cardmethod \\
    LA object     & \cardobject \\
    Algorithms    & \cardalgorithms \\
    Data          & \carddata \\
    \ifnum\stage>1 Code & \cardcode \\ \fi
    \ifnum\stage=3 Recording & \cardvideo \\ \fi
    \bottomrule
  \end{tabularx}\par}

% Page counting: the main body runs from \bodystart to \bodyend
\newcommand{\pagelimit}{\ifnum\stage=3 10\else 2\fi}
\newcommand{\bodystart}{\label{body-start}}
\newcommand{\bodyend}{\label{body-end}}
\AtEndDocument{%
  \refused{body-start}\refused{body-end}%
  \edef\bodypages{\number\numexpr\getpagerefnumber{body-end}
    -\getpagerefnumber{body-start}+1\relax}%
  \typeout{}%
  \typeout{*** Main body: \bodypages\space page(s); limit \pagelimit.}%
  \typeout{}%
  \ifnum\bodypages>\pagelimit\relax
    \PackageWarningNoLine{LA-template}{Main body has \bodypages\space pages,
      above the limit of \pagelimit}%
  \fi}

%%%%%%%%%%%%%%%%%%%%%%%%%%%%%%%%%%%%%%%%%%%%%%%%%%%%%%%%%%%%%%%%%%%%%%%%%
\begin{document}
%%%%%%%%%%%%%%%%%%%%%%%%%%%%%%%%%%%%%%%%%%%%%%%%%%%%%%%%%%%%%%%%%%%%%%%%%

\ifnum\stage=3
%----------------------------------------------------------------------
%  Final report: title page (not counted)
%----------------------------------------------------------------------
\begin{titlepage}
  \centering
  {\large\scshape Ukrainian Catholic University\par}
  {\scshape Faculty of Applied Sciences\quad$\cdot$\quad
    Data Science Master Programme\par}
  \vspace{1.6cm}
  \rule{\linewidth}{0.5mm}\par\vspace{0.4cm}
  {\LARGE\bfseries \reporttitle\par}
  \vspace{0.3cm}
  {\large Linear Algebra: course project, final report\par}
  \vspace{0.2cm}
  \rule{\linewidth}{0.5mm}\par
  \vspace{1.2cm}
  {\large\emph{Authors:}\par\smallskip \reportauthors\par}
  \vspace{1.2cm}
  \projectcard
  \vfill
    \IfFileExists{UCU-Apps.png}{\includegraphics[height=4cm]{UCU-Apps.png}\par}{}
  \vspace{0.2cm}
  {\large \today\par}
\end{titlepage}

%----------------------------------------------------------------------
%  Abstract: opens the first page of the main body
%----------------------------------------------------------------------
\bodystart
\begin{abstract}
\guide{Five to eight sentences: the problem, the linear-algebra object,
the methods you compared, the main quantitative result, and the main
limitation. Write it last.}
We compare two algorithms for completing a partially observed ratings
matrix by a matrix of low rank~\cite[Sec.~III.5]{Strang}.
\end{abstract}

\else
%----------------------------------------------------------------------
%  Interim reports: compact header and project card (counted)
%----------------------------------------------------------------------
\bodystart
\begin{center}
  {\Large\bfseries \reporttitle\par}
  \medskip
  {\large \stagename\par}
  \smallskip
  \begin{tabular}{c}\reportauthors\end{tabular}\par
  \smallskip
  {\small Linear Algebra, UCU Data Science Master Programme, \today\par}
\end{center}
\projectcard
\medskip
\fi


%%%%%%%%%%%%%%%%%%%%%%%%%%%%%%%%%%%%%%%%%%%%%%%%%%%%%%%%%%%%%%%%%%%%%%%%%
\section{Problem and motivation}
\label{sec:problem}

The detection of AI-generated music has proven a problem
for platforms, rights holders and listeners. In this work,
we intend to distinguish AI-generated audio from human-generated audio.
Given a short audio clip, we predict human-origin versus text-to-music-generated.
We do so by checking whether the low-rank structure of its STFT magnitude
spectrogram supports this distinction on held-out clips,
rather than create a universal detector. Prior work also warns that
strong controlled results may not generalize across manipulations
or datasets~\cite{Afchar,LiBench}, so we try to balance classes
and group each original/synthetic caption pair in the same data split.

Success is held-out balanced accuracy $\geq0.65$, with the lower
bound of a 95\% confidence interval above $0.50$.
Separately, randomized SVD should achieve relative $\ell_2$
error $\leq10^{-2}$ in the first 20 normalized squared singular-value
features on at least 95\% of test clips, relative to exact SVD.
If the classification threshold is not met,
we should deem these features insufficient in this setup.

%%%%%%%%%%%%%%%%%%%%%%%%%%%%%%%%%%%%%%%%%%%%%%%%%%%%%%%%%%%%%%%%%%%%%%%%%
\section{Linear-algebra formulation}
\label{sec:formulation}

Each clip is resampled to 22.05 kHz. A centred STFT
(window/FFT length 2048, hop 512) gives a nonnegative
magnitude matrix $M\in\mathbb{R}_{\geq0}^{F\times T}$,
where $F=1025$ and $T\approx431$ for ten seconds.
Its singular value decomposition is

\begin{equation}\label{eq:svd}
  M=U\Sigma V^\top,\qquad \sigma_1\geq\sigma_2\geq\cdots\geq0.
\end{equation}

We use $z=(p_1,\ldots,p_{20})$ as features for logistic regression,
where $p_i=\sigma_i^2/\|M\|_F^2$ and $\|M\|_F^2=\sum_i\sigma_i^2$.
The features summarize the distribution of spectrogram energy across
orthogonal components.

The LA task is to compute or approximate the leading singular values
and test whether this low-rank summary preserves class information.
Because the features discard phase and temporal ordering,
useful cues may be lost even with an accurate decomposition.

%%%%%%%%%%%%%%%%%%%%%%%%%%%%%%%%%%%%%%%%%%%%%%%%%%%%%%%%%%%%%%%%%%%%%%%%%
\section{Methods}
\label{sec:methods}

\subsection{Exact thin SVD}

A standard dense SVD gives the comparison values and
the optimal rank-$k$ approximation for the Frobenius norm
by the Eckart--Young--Mirsky theorem.

For $F\geq T$, its cost is $O(FT^2+T^3)$ flops and $O(FT)$ storage.
It is accurate and straightforward, but repeats work for every clip.
In our case, we will have moderate matrix sizes, and so it may already be fast,
limiting the practical speedup available to a randomized method.
Small singular values can also be sensitive to rounding,
and an accurate decomposition does not guarantee useful detection features.

NumPy`s \textrm{np.linalg.svd} uses LAPACK`s \textrm{\_gesdd} routine
to calculate SVD. \textrm{\_gesdd} implements a divide-and-conquer algorithm
on a bidiagonal (diagonal with entries non-zero above or below main diagonal)
matrix. It divides the matrix into sizes that are convenient for QR iteration
for SVD calculation, and then accumulates the results.

We left out the algorithm pseudo-code for the \textrm{\_gesdd} implementation for brevity.

% \begin{algorithm}[ht]
% \caption{Householder divide-and-conquer thin SVD}\label{alg:householder_dnc_svd}
%
% \KwIn{$A\in\mathbb{R}^{F\times T}$ with $F\ge T$}
%
% \KwOut{Thin SVD factors: $U\in\mathbb{R}^{F\times T}$, singular values $\mathbf{s}\in\mathbb{R}^{T}$, $V\in\mathbb{R}^{T\times T}$}
%
% \tcp*{Phase 1: Householder reduction to upper bidiagonal form}
%
% Compute Householder reflections to reduce $A$ to bidiagonal $B\in\mathbb{R}^{T\times T}$:
%
% $A \gets U_1 B V_1^\top$\tcp*{Cost: $O(FT^2 - \frac{1}{3}T^3)$ flops}
%
% \medskip
%
% \SetKwProg{Fn}{Function}{}{end}\SetKwFunction{FRecurs}{FnRecursive}%
% \newcommand{\forcond}{$i=0$ \KwTo $n$}
%
% \tcp*{Phase 2: Divide-and-conquer SVD on the bidiagonal matrix}
% \Fn{\text{BidiagDivideAndConquer}($B$)}{
% \If{$\text{size}(B)$ is small enough}{
% Compute SVD directly via QR iteration: $B = \tilde{U}\operatorname{diag}(\mathbf{s})\tilde{V}^\top$;
% \Return $\tilde{U}$, $\mathbf{s}$, $\tilde{V}$;
% }
%
% Split $B$ into independent sub-bidiagonal blocks $B_1$ and $B_2$;
% $(\tilde{U}_1, \mathbf{s}_1, \tilde{V}_1) \gets \text{BidiagDivideAndConquer}({B}_1)$;
% $(\tilde{U}_2, \mathbf{s}_2, \tilde{V}_2) \gets \text{BidiagDivideAndConquer}({B}_2)$;
%
% Merge subproblems using secular equations to obtain $B = U_2 \operatorname{diag}(\mathbf{s}) V_2^\top$;
%
% \Return $U_2$, $\mathbf{s}$, $V_2$;
% }
%
% $(U_2, \mathbf{s}, V_2) \gets \text{BidiagDivideAndConquer}(B)$;
% \medskip
%
% \tcp*{Phase 3: Accumulate orthogonal transformations}
%
% $U \gets U_1 U_2$;\quad $V \gets V_1 V_2$;
%
% \Return{$U$, $\mathbf{s}$, $V$}
% \end{algorithm}

\medskip\noindent\textbf{Cost.} Phase 1 requires $O(FT^2 - \frac{1}{3}T^3)$
flops for Householder bidiagonalization.
Phase 2 (divide-and-conquer) and Phase 3 (vector accumulation) require $O(T^3)$ flops.
Total arithmetic cost is $O(FT^2 + T^3)$, but the divide-and-conquer conquering
step heavily benefits from parallel matrix-matrix operations,
making it much faster in practice than sequential QR iteration for large dimensions.

\begin{lstlisting}[caption={Singular Value Decomposition in NumPy.},label={lst:power}]
import numpy as np

def householder_dnc_thin_svd(A):
    U, s, Vt = np.linalg.svd(A)
    return U, s, Vt
\end{lstlisting}

\subsection{Randomized SVD}

Following the randomized range-finding framework of Halko, Martinsson,
and Tropp~\cite{Halko}, the core idea of randomized SVD is to capture the essential column space (range)
of a matrix $A \in \mathbb{R}^{F \times T}$ using a low-dimensional random projection before
performing an exact factorization. The procedure unfolds in four main stages:

We wish to find a low-rank approximation of rank $k$.
Instead of examining all $T$ dimensions, we draw a random test matrix
$\mathbf{\Omega} \in \mathbb{R}^{T \times \ell}$ whose entries are independent
and identically distributed standard Gaussian variables, $\mathbf{\Omega}_{ij} \sim \mathcal{N}(0, 1)$.

The sketch dimension $\ell$ is chosen slightly larger than
the target rank $k$ via an oversampling parameter $p$
(e.g., $k=20$, $p=10$, yielding $\ell = 30$).

Multiplying $A$ by $\mathbf{\Omega}$ forms the sketch matrix:
\begin{equation}
  Y = A\mathbf{\Omega} \in \mathbb{R}^{F \times \ell}
\end{equation}

If the singular values of $A$ decay slowly, the random sketch
alone may capture too much noise from the tail of the spectrum.
To sharpen the separation between the dominant components and the
trailing noise, we apply $q=1$ power iteration steps (subspace iteration).
This alternates multiplications by $A$ and $A^\top$ to damp out lesser singular values:
\begin{align}
  \tilde{Y} &= A^\top Q, \quad &\text{QR factorization: } \tilde{Y} = \tilde{Q} \tilde{R} \\
  Y &= A \tilde{Q}, \quad &\text{QR factorization: } Y = Q R
\end{align}

This accelerates the spectral decay rate without requiring a full matrix factorization.

Once the refined orthonormal basis $Q$ is established,
we project the original matrix $A$ onto this lower-dimensional subspace to form a much smaller matrix $B$:

\begin{equation}
  B = Q^\top A \in \mathbb{R}^{\ell \times T}
\end{equation}

Because $\ell \ll T$, computing the exact SVD of $B$ via
standard routines (e.g., LAPACK's $\_gesdd$) is relatively fast.

\begin{equation}
  B = \tilde{U} \mathbf{s} V^\top
\end{equation}

Finally, we map the left singular vectors back to the original
space by setting $U = Q \tilde{U}$, and truncate all factors to the target rank $k$.

\begin{algorithm}[ht]
  \caption{Randomized SVD with Power Step}\label{alg:rsvd}

  \KwIn{Matrix $A \in \mathbb{R}^{F \times T}$ ($F \ge T$), target rank $k$, oversampling parameter $p$, power iterations $q$}

  \KwOut{Rank-$k$ SVD factors: $U_k \in \mathbb{R}^{F \times k}$, $\mathbf{s}_k \in \mathbb{R}^{k}$, $V_k \in \mathbb{R}^{T \times k}$}

  $\ell \gets k + p$\tcp*{Total sketch dimension, $\ell \approx 30$}

  Draw Gaussian random matrix $\mathbf{\Omega} \in \mathbb{R}^{T \times \ell}$ where $\mathbf{\Omega}_{ij} \sim \mathcal{N}(0, 1)$\;

  $Y \gets A \mathbf{\Omega}$\tcp*{Form initial sketch: $O(FT\ell)$ flops}

  \tcp*{Power iteration to accelerate singular value decay}
  \For{$i = 1 \text{ to } q$}{
    $Q, \_ \gets \operatorname{qr}(Y)$\;
    $\tilde{Y} \gets A^\top Q$\;
    $\tilde{Q}, \_ \gets \operatorname{qr}(\tilde{Y})$\;
    $Y \gets A \tilde{Q}$\;
  }
  $Q, \_ \gets \operatorname{qr}(Y)$\tcp*{Orthonormal basis for range of $A$: $O(F\ell^2)$ flops}

  \medskip
  \tcp*{Project and perform small SVD}
  $B \gets Q^\top A$\tcp*{Compress matrix to dimension $\ell \times T$: $O(FT\ell)$ flops}
  $\tilde{U}, \mathbf{s}, V^\top \gets \operatorname{SVD}(B)$\tcp*{Exact SVD on reduced matrix: $O(T\ell^2)$ flops}

  \medskip
  \tcp*{Reconstruct full left singular vectors and truncate to rank $k$}
  $U \gets Q \tilde{U}$\;
  $U_k \gets U[:, 1:k]$;\quad $\mathbf{s}_k \gets \mathbf{s}[1:k]$;\quad $V_k \gets V[:, 1:k]$\;
  \Return{$U_k$, $\mathbf{s}_k$, $V_k$}
\end{algorithm}

\medskip\noindent\textbf{Cost.}
With $q$ power steps, the overall flop count is $O((2q+1)FT\ell + (F+T)\ell^2)$,
and memory requirement is bounded by $O(FT + (F+T)\ell)$. Because $\ell \ll T$,
the dense SVD phase runs in $O(T\ell^2)$ time instead of $O(T^3)$.

\begin{lstlisting}[caption={Randomized SVD with $q=1$ power step in NumPy.},label={lst:rsvd}]
import numpy as np

def randomized_svd(A, k=20, p=10, q=1):
    F, T = A.shape
    ell = k + p

    # Draw Gaussian random test matrix
    Omega = np.random.normal(size=(T, ell))
    Y = A @ Omega

    # Power iteration (q steps)
    for _ in range(q):
        Q, _ = np.linalg.qr(Y)
        Y_tilde = A.T @ Q
        Q_tilde, _ = np.linalg.qr(Y_tilde)
        Y = A @ Q_tilde

    # Final orthogonalization for the range space
    Q, _ = np.linalg.qr(Y)

    # Project A into smaller subspace and compute exact SVD
    B = Q.T @ A
    U_tilde, s, Vt = np.linalg.svd(B, full_matrices=False)

    # Reconstruct original left singular vectors
    U = Q @ U_tilde

    # Truncate to target rank k
    return U[:, :k], s[:k], Vt[:k, :].T
\end{lstlisting}

\subsection{Cost and failure modes}

Randomized SVD is a low-rank approximation of thin SVD.
It only computes the singular value for a random subspace of a matrix,
and might be vulnerable to flat spectra (ones, where $|\sigma_{k+1}/\sigma_k| \approx 1$).
Unlike regular SVD, randomized SVD does not rely on reduction to the bidiagonal form,
rather, the bulk of the computation is the matrix multiplication during the power iteration step.

Exact SVD requires storing intermediate orthogonal transformation
matrices and workspaces that scale with $T^2$ and randomized SVD operates primarily on
compressed sketch matrices of size $\ell = k + p$ (where $\ell \ll T$), and so it requires substantially less memory.

Randomized SVD does not handle ill-conditioned matrices well.
Should the spectrum of a matrix decay slowly, the power iteration step
would not capture the matrice`s spectrum in the randomized subspace well.

%%%%%%%%%%%%%%%%%%%%%%%%%%%%%%%%%%%%%%%%%%%%%%%%%%%%%%%%%%%%%%%%%%%%%%%%%
\section{Implementation and experiments}
\label{sec:experiments}


\subsection{Setup}

The project will be implemented in Python using NumPy and SciPy for
linear algebra and signal processing, with scikit-learn for the
baseline classifier and evaluation metrics.

Audio preprocessing should use a common pipeline:
conversion to mono at 16\,kHz, low-pass
filtering at 7\,kHz, loudness normalisation, and a common
MP3 encode/decode step.

The datasets are FakeMusicCaps and the corresponding human-made MusicCaps
clips, paired by their original captions. We will
balance the classes and split the data by caption
to prevent related examples from appearing in both training and
test sets.

For each clip, we need to compute a log-magnitude STFT spectrogram and
extract singular-value features, including entropy-based effective
rank and energy concentration. We also compute principal angles
between the leading left singular subspaces of consecutive
spectrogram windows. A simple bandwidth-only classifier will serve
as a diagnostic baseline for detecting pipeline-related differences.
The current implementation is being developed and tested on small
subsets before the full dataset is processed.

\subsection{Results}

The initial experiments will compare the proposed features with the
bandwidth-only baseline using balanced accuracy and ROC-AUC on a
held-out test set. We will also investigate the relationship between
the singular-value spectrum and temporal subspace stability, and
whether these features generalise across different generated clips.

Exact SVD is expected to be inexpensive for individual
spectrograms, so it will be used as the reference method for
per-clip features. Randomized SVD will instead be applied to a
larger matrix formed by stacking spectrogram frames from training
clips, allowing us to learn a shared spectral basis.

We will report quantitative results on real audio
after we retrieve and process the datasets.

%%%%%%%%%%%%%%%%%%%%%%%%%%%%%%%%%%%%%%%%%%%%%%%%%%%%%%%%%%%%%%%%%%%%%%%%%
\section{Discussion and limitations}
\label{sec:discussion}

The main risk is that AI-generated and human-made clips may come
from different audio pipelines. Differences in sampling rate,
YouTube compression, bandwidth, mastering, or loudness could create
spectral cues unrelated to the process that generated the music.
Resampling, common low-pass filtering, loudness normalisation, and
common codec processing should reduce some of these differences,
but cannot guarantee their removal. The bandwidth-only baseline
will help reveal whether simple technical artifacts already
predict the class labels.

A second limitation is that low-rank spectral structure is not
necessarily unique to AI-generated music. Human recordings and
generated music can both contain strong redundancy, while different
generators may produce different spectral characteristics.
For this, principal angles between consecutive windows should provide
a concrete measure of subspace change,
but, for classification, they might not be of much use.

%%%%%%%%%%%%%%%%%%%%%%%%%%%%%%%%%%%%%%%%%%%%%%%%%%%%%%%%%%%%%%%%%%%%%%%%%
\section{Conclusion}
\label{sec:conclusion}

\guide{What you achieved, the main difficulties, and what you would do
next. Half a page is enough.}

\ifnum\stage<3
%%%%%%%%%%%%%%%%%%%%%%%%%%%%%%%%%%%%%%%%%%%%%%%%%%%%%%%%%%%%%%%%%%%%%%%%%
\section{Plan and division of work}
\label{sec:plan}

\guide{Interim reports only; this section disappears from the final
report. \emph{Interim report 1}: the planned steps and who does what.
\emph{Interim report 2}: what is not working yet, any change of plan,
and the outline of the final report.}
\fi

\bodyend

%%%%%%%%%%%%%%%%%%%%%%%%%%%%%%%%%%%%%%%%%%%%%%%%%%%%%%%%%%%%%%%%%%%%%%%%%
%%  References (not counted)
%%%%%%%%%%%%%%%%%%%%%%%%%%%%%%%%%%%%%%%%%%%%%%%%%%%%%%%%%%%%%%%%%%%%%%%%%
% Cite only sources you have actually read. Give a DOI or a link where
% one exists. Cite the datasets and the main software you used.
\begin{thebibliography}{9}

\bibitem{FakeMusicCaps}
L. Comanducci, P. Bestagini and S. Tubaro,
``FakeMusicCaps: A Dataset for Detection and Attribution of Synthetic Music Generated via Text-to-Music Models,''
\emph{Journal of Imaging}, vol. 11, no. 7, article 242, 2025.
\href{https://doi.org/10.3390/jimaging11070242}{doi:10.3390/jimaging11070242}.

\bibitem{Afchar}
D. Afchar, G. Meseguer-Brocal and R. Hennequin,
``AI-Generated Music Detection and its Challenges,'' in
\emph{Proc. IEEE International Conference on Acoustics, Speech and Signal Processing (ICASSP)}, 2025.
\href{https://doi.org/10.1109/ICASSP49660.2025.10890655}{doi:10.1109/ICASSP49660.2025.10890655}.

\bibitem{LiBench}
Y. Li, Q. Sun, H. Li, L. Specia and B. W. Schuller,
``Explainable detection of machine generated music and early systematic evaluation,''
\emph{Scientific Reports}, vol. 16, article 13757, 2026.
\href{https://doi.org/10.1038/s41598-026-42133-7}{doi:10.1038/s41598-026-42133-7}.

\bibitem{Halko}
N. Halko, P. G. Martinsson and J. A. Tropp,
``Finding Structure with Randomness: Probabilistic Algorithms for Constructing Approximate Matrix Decompositions,''
\emph{SIAM Review}, vol. 53, no. 2, pp. 217--288, 2011.
\href{https://doi.org/10.1137/090771806}{doi:10.1137/090771806}.

\end{thebibliography}

%%%%%%%%%%%%%%%%%%%%%%%%%%%%%%%%%%%%%%%%%%%%%%%%%%%%%%%%%%%%%%%%%%%%%%%%%
%%  Appendices (not counted)
%%%%%%%%%%%%%%%%%%%%%%%%%%%%%%%%%%%%%%%%%%%%%%%%%%%%%%%%%%%%%%%%%%%%%%%%%
\ifnum\stage=3
\appendix

\section{Contribution statement}

\guide{Required. Who did what, agreed by all members. Staff may adjust
an individual grade if contributions are clearly unequal.}

\noindent
\begin{tabularx}{\linewidth}{@{}lX@{}}
  \toprule
  Member & Contribution \\
  \midrule
  A.~A.~Lovelace & \dots \\
  C.~Babbage     & \dots \\
  \bottomrule
\end{tabularx}

\section{AI-use statement}

\guide{Required. Which AI tools you used and for what: code, language
editing, finding literature. Every member must be able to explain every
equation and every line of the core code. If you used none, say so.}

\section{Additional results}

\guide{Optional. Proofs, extra figures and tables that support the main
text. Reviewers are not required to read the appendices, so everything
needed to follow the report belongs in the main body.}
\fi

\end{document}
